\clearpage
\section{Limits}

Limit proofs hinge on the following definition of a limit:

\[\lim_{x \to a} f(a) = L \iff \forall \varepsilon > 0, \]
\[\exists\ \delta\ s.t.\ 0 < |x - a| < \delta \implies |f(x) - L| < \varepsilon\]


\subsection{Addition Rule}

Statement:

given that 
\[\lim_{x \to a} f(x) = L, \lim_{x \to a} f(x) = M\], show that
\[\lim_{x \to a} f(x) + g(x) = L + M\]

\begin{proof}
Consider some $\varepsilon > 0$, then by the definition of the limit there exists $\delta_1$ and $\delta_2$ such that:

\[0 < |x - a| < \delta_1 \implies |f(x) - L| < \frac{\varepsilon}{2}\]

\[0 < |x - a| < \delta_2 \implies |g(x) - M| < \frac{\varepsilon}{2}\]

Then let $\delta = \text{min}(\delta_1, \delta_2)$, and we need
\[|f(x) + g(x) - (L + M)| < \varepsilon\]

We have:

\[0 < |x - a| < \delta\]
which implies
\[|f(x) - L| < \frac{\varepsilon}{2}\]
and
\[|g(x) - M| < \frac{\varepsilon}{2}\]

By the triangle inequality, 
\[|a + b| \leq |a| + |b|\]
so
\[|f(x) + g(x) - (L + M)| < |f(x) - L| + |g(x) - L| < \frac{\varepsilon}{2} + \frac{\varepsilon}{2} = \varepsilon\]
so
\[|f(x) + g(x) - (L + M)| < \varepsilon\]
and so limits follow the addition rule.
\qed
\end{proof}

Do note that I \emph{did} need to consult the textbook for the triangle inequality stuff, but the rest was fine.

\subsection{Multiplication Rule}
Given that
\[\lim_{x \to a} f(x) = L\]
and
\[\lim_{x \to a} g(x) = M\]
prove that
\[\lim_{x \to a} f(x)g(x) = LM\]

\begin{proof}

\end{proof}


\subsection{Polynomials are continuous}

A function, $f(x)$, is continuous if \[\lim_{x \to a} f(x) = f(a)\]

From this, let us prove that all polynomials of the form $x^n$ for $n \in \mathbb{N}$ are continuous. 
